%!TEX program = xelatex
% Compilar: xelatex hasse_diagram_report.tex  (dos pasadas para tikz)
\documentclass[11pt, a4paper]{article}

% ── Fuentes (XeLaTeX) ─────────────────────────────────────────────────────────
\usepackage{fontspec}
\setmainfont{Segoe UI}[
  BoldFont       = Segoe UI Bold,
  ItalicFont     = Segoe UI Italic,
  BoldItalicFont = Segoe UI Bold Italic
]
\setsansfont{Segoe UI}[
  BoldFont       = Segoe UI Bold,
  ItalicFont     = Segoe UI Italic,
  BoldItalicFont = Segoe UI Bold Italic
]
\setmonofont{Consolas}[
  BoldFont = Consolas Bold
]

% ── Geometría ─────────────────────────────────────────────────────────────────
\usepackage[
  a4paper,
  top         = 2.8cm,
  bottom      = 2.4cm,
  left        = 2.2cm,
  right       = 2.2cm,
  headheight  = 30pt,
  headsep     = 14pt
]{geometry}

% ── Colores ───────────────────────────────────────────────────────────────────
\usepackage[table]{xcolor}
\definecolor{cNavy}   {HTML}{0D1B3E}
\definecolor{cBlue}   {HTML}{2A4F9B}
\definecolor{cAccent} {HTML}{4A7FD4}
\definecolor{cBody}   {HTML}{1A1A2E}
\definecolor{cMuted}  {HTML}{6B7FA8}
\definecolor{cCodeFG} {HTML}{1B3A70}
\definecolor{cCodeBG} {HTML}{EEF2FA}
\definecolor{cTHbg}   {HTML}{0D1B3E}
\definecolor{cTRalt}  {HTML}{F2F5FC}
\definecolor{cRule}   {HTML}{C0CADF}
\definecolor{cQuoteBG}{HTML}{F0F5FF}
\definecolor{cGold}   {HTML}{B8860B}

% ── Matemáticas ───────────────────────────────────────────────────────────────
\usepackage{amsmath}
\usepackage{amssymb}
\usepackage{mathtools}

% ── Tipografía ────────────────────────────────────────────────────────────────
\usepackage{microtype}
\usepackage{parskip}
\setlength{\parindent}{0pt}
\setlength{\parskip}{5pt plus 1pt}

% ── Tablas ────────────────────────────────────────────────────────────────────
\usepackage{booktabs}
\usepackage{tabularx}
\usepackage{colortbl}
\usepackage{array}
\usepackage{adjustbox}
\renewcommand{\arraystretch}{1.4}
\newcolumntype{L}[1]{>{\raggedright\arraybackslash}p{#1}}
\newcolumntype{C}[1]{>{\centering\arraybackslash}p{#1}}

% ── Cajas ─────────────────────────────────────────────────────────────────────
\usepackage{tcolorbox}
\tcbuselibrary{skins, breakable}

\tcbset{
  codebox/.style = {
    enhanced,
    colback    = cCodeBG,
    colframe   = cAccent,
    arc        = 3pt,
    leftrule   = 3pt,
    rightrule  = 0.5pt,
    toprule    = 0.5pt,
    bottomrule = 0.5pt,
    left       = 10pt,
    right      = 10pt,
    top        = 7pt,
    bottom     = 7pt,
    fontupper  = \small\ttfamily\color{cCodeFG},
    breakable
  },
  quotebox/.style = {
    enhanced,
    colback    = cQuoteBG,
    colframe   = cAccent,
    arc        = 2pt,
    leftrule   = 3pt,
    rightrule  = 0pt,
    toprule    = 0pt,
    bottomrule = 0pt,
    left       = 12pt,
    right      = 10pt,
    top        = 7pt,
    bottom     = 7pt,
    fontupper  = \itshape\color{cCodeFG},
    breakable
  },
  theorembox/.style = {
    enhanced,
    colback    = cCodeBG,
    colframe   = cGold,
    arc        = 3pt,
    leftrule   = 4pt,
    right      = 10pt,
    left       = 12pt,
    top        = 8pt,
    bottom     = 8pt,
    breakable
  }
}

% ── Encabezados de sección ────────────────────────────────────────────────────
\usepackage{titlesec}

% Secciones numeradas con contador personalizado (§ 0, § 1, ...)
\titleformat{\section}[block]
  {\color{cNavy}\Large\bfseries}
  {\color{cAccent}\S\,\thesection}{0.6em}{}
  [\vspace{2pt}{\color{cAccent}\hrule height 1.6pt}\vspace{4pt}]

\titleformat{\subsection}
  {\color{cBlue}\large\bfseries}
  {}{0em}{}
  [\vspace{1pt}{\color{cRule}\hrule height 0.5pt}\vspace{2pt}]

\titleformat{\subsubsection}
  {\color{cBlue}\normalsize\bfseries}
  {}{0em}{}

\titlespacing*{\section}      {0pt}{20pt}{8pt}
\titlespacing*{\subsection}   {0pt}{14pt}{4pt}
\titlespacing*{\subsubsection}{0pt}{10pt}{3pt}

% Permitir sección 0
\setcounter{section}{-1}

% ── Cabecera y pie de página ──────────────────────────────────────────────────
\usepackage{fancyhdr}
\usepackage{tikz}
\usetikzlibrary{calc}

\pagestyle{fancy}
\fancyhf{}
\renewcommand{\headrulewidth}{0pt}

\fancyhead[C]{%
  \begin{tikzpicture}[remember picture, overlay]
    \fill[cNavy]
      (current page.north west)
      rectangle ($(current page.north east) + (0, -1.15cm)$);
    \fill[cAccent]
      (current page.north west)
      rectangle ($(current page.north west) + (0.28\paperwidth, -1.15cm)$);
    \node[white, anchor=west, font=\small\itshape] at
      ($(current page.north west) + (0.28\paperwidth + 8pt, -0.58cm)$)
      {Go Ising \enspace\textbullet\enspace Experimento 06
       \enspace\textbullet\enspace Diagrama de Hasse · Orden Parcial Pareto};
  \end{tikzpicture}%
}

\fancyfoot[L]{\small\color{cMuted} Leonardo Jiménez Martínez \enspace·\enspace UNAM \enspace·\enspace 2026}
\fancyfoot[R]{\small\color{cMuted} Página \thepage}
\def\footrule{{\color{cRule}\vspace{-3pt}\hrule width\headwidth height 0.4pt\vspace{2pt}}}

% ── Listas ────────────────────────────────────────────────────────────────────
\usepackage{enumitem}
\setlist[itemize]{
  leftmargin = 1.5em,
  itemsep    = 2pt,
  topsep     = 4pt,
  label      = {\color{cAccent}\small$\bullet$}
}
\setlist[enumerate]{leftmargin=1.8em, itemsep=3pt, topsep=4pt}

% ── PDF metadata ──────────────────────────────────────────────────────────────
\usepackage[
  colorlinks = true,
  linkcolor  = cAccent,
  urlcolor   = cAccent,
  pdfauthor  = {Leonardo Jiménez Martínez},
  pdftitle   = {Diagrama de Hasse del Orden Parcial Pareto},
  pdfsubject = {Go Ising · Experimento 06 · Hamiltonians cúbicos}
]{hyperref}

% ── Macros ────────────────────────────────────────────────────────────────────
\newcommand{\TH}[1]{\cellcolor{cTHbg}\color{white}\textbf{#1}}
\newcommand{\frente}[1]{\textcolor{cAccent}{\textbf{F#1}}}

% =============================================================================
\begin{document}
\thispagestyle{fancy}

% ── Portada ───────────────────────────────────────────────────────────────────
\vspace*{4pt}
{\color{cNavy}\Huge\bfseries Diagrama de Hasse del Orden Parcial Pareto}

\vspace{6pt}
{\color{cMuted}\large\itshape
  Interpretación matemática y estratégica para el juego de Go}

\vspace{4pt}
{\color{cMuted}\small\itshape
  Experimento 06 --- Familia de Hamiltonianos cúbicos}\\[1pt]
{\color{cMuted}\small\itshape
  296 candidatos \enspace·\enspace
  148 frentes \enspace·\enspace
  475 relaciones de cobertura}

\vspace{8pt}
{\color{cAccent}\hrule height 2pt}
\vspace{10pt}

% =============================================================================
\section{La familia de candidatos: ¿qué son estos Hamiltonianos?}

\subsection{Go como sistema de espines}

Cada intersección del tablero de Go tiene un valor de espín
$s \in \{-1, 0, +1\}$ (negro, vacío, blanco). La energía de interacción entre
dos intersecciones vecinas $i$ y $j$ está dada por un \textbf{Hamiltoniano de
par} $H(s_i, s_j)$. La energía total del tablero es la suma sobre todos los
pares adyacentes:

\begin{tcolorbox}[codebox]
$\displaystyle E \;=\; \sum_{\langle i,j \rangle} H(s_i,\; s_j)$
\end{tcolorbox}

El modelo de evaluación estratégica se reduce completamente a la elección de
$H$: qué tipo de vecindades favorece o penaliza.

\subsection{Por qué polinomios cúbicos}

Sobre $\{-1, 0, +1\}$, se cumple $x^3 = x$ para todo $x$. Por eso los monomios
de grado $\geq 3$ en una sola variable no aportan información nueva. Los
\textbf{monomios independientes} en $(x,y) \in \{-1,0,+1\}^2$ son exactamente
7 (sin contar la constante):

\begin{tcolorbox}[codebox]
$x, \quad y, \quad x^2, \quad xy, \quad y^2, \quad x^2 y, \quad xy^2$
\end{tcolorbox}

Cualquier función $H\colon\{-1,0,+1\}^2\to\mathbb{R}$ se expresa como
combinación lineal de estos 7 términos. Esto hace que los polinomios cúbicos
sean la \textbf{familia completa} --- no existe una familia más general para
este espacio de estados.

\subsection{La plantilla \texttt{cubic\_mixed}}

El experimento exploró la plantilla:

\begin{tcolorbox}[codebox]
$H(x,y) \;=\; a_1 x + a_2 y + b_{11} x^2 + b_{12}\,xy + b_{22} y^2
              + c_{112}\,x^2 y + c_{122}\,x y^2$
\end{tcolorbox}

Los coeficientes se muestrearon aleatoriamente (\texttt{seed = 42},
300 muestras) con:

\begin{itemize}
  \item $a_1,\; a_2,\; b_{11},\; b_{12},\; b_{22} \;\in\; [-3,\; 3]$
  \item $c_{112},\; c_{122} \;\in\; [-1,\; 1]$
\end{itemize}

Cada nodo del diagrama de Hasse representa uno de los \textbf{296 candidatos}
que pasaron el filtro de calidad (de los 305 analizados, incluyendo las
referencias: $H_{M1}$, Alvarado y el cúbico armónico).

\subsection{Qué mide cada criterio de calidad}

Los 4 criterios del orden Pareto evalúan propiedades distintas e independientes:

\medskip
\begin{center}
\rowcolors{2}{white}{cTRalt}
\begin{tabularx}{0.96\textwidth}{L{2.5cm}L{4cm}X}
  \TH{Criterio} & \TH{Cómo se calcula} & \TH{Qué mide para Go} \\[1pt]
  $H_1^{\max}$
    & Vida máxima de un generador $H_1$ en la filtración por subniveles de
      $H$ sobre $[-2,2]^2$, normalizada por el rango
    & Capacidad del Hamiltoniano de detectar \textbf{estructuras cíclicas de
      influencia}: grupos que rodean a otros, cadenas de conexión, seki \\
  Robustez
    & Fracción de 12 perturbaciones $\pm 5\%$ de los coeficientes que
      mantienen $H_1^{\max} > 0.20$
    & \textbf{Estabilidad estructural}: las propiedades del modelo no son
      artefactos de coeficientes exactos; funcionarán bajo variación numérica \\
  $\Delta E$
    & $H(x^*, y^*)_{\max} - H(x^*, y^*)_{\min}$ sobre $[-2,2]^2$
    & \textbf{Contraste energético}: separación entre posiciones de máxima
      y mínima energía; mayor $\Delta E$ = mayor poder discriminativo \\
  $n_{A_1}$
    & Número de puntos críticos con $\det(\mathrm{Hess}\,H) < 0$ (sillas)
    & \textbf{Transiciones de régimen}: cada nodo $A_1$ es un punto donde la
      topología de la fibra $H^{-1}(c)$ cambia cualitativamente \\
\end{tabularx}
\end{center}
\medskip

Los dos primeros criterios ($H_1^{\max}$, Robustez) miden riqueza y estabilidad
topológica. Los dos últimos ($\Delta E$, $n_{A_1}$) miden potencia discriminativa
y complejidad estructural. Los 4 son genuinamente independientes: ningún candidato
maximiza los 4 simultáneamente---de ahí que el poset no tenga supremo.

% =============================================================================
\section{Qué es este diagrama}

El diagrama de Hasse es la representación visual de un \textbf{orden parcial}:
una forma de comparar candidatos que no obliga a elegir un único ``mejor'' sino
que admite múltiples óptimos incomparables entre sí.

Cada nodo es un \textbf{Hamiltoniano polinómico cúbico} $H(x,y)$ que asigna
una energía a cada par de intersecciones de Go. La imagen dentro de cada nodo
muestra la variedad $\Gamma(H) = \{z = H(x,y)\} \subset \mathbb{R}^3$ --- la
``forma'' topológica de ese Hamiltoniano.

% =============================================================================
\section{La relación de orden: qué significa ``dominar''}

Cada candidato $H$ tiene cuatro valores que miden su calidad como función de
evaluación estratégica:

\medskip
\begin{center}
\rowcolors{2}{white}{cTRalt}
\begin{tabularx}{0.96\textwidth}{L{2.4cm}C{2cm}X}
  \TH{Criterio} & \TH{Símbolo} & \TH{Qué mide} \\[1pt]
  Vida máxima $H_1$
    & $H_1^{\max}$
    & Ciclos topológicos en la fibración de Milnor --- complejidad táctica \\
  Robustez
    & Rob
    & Fracción de perturbaciones $\pm 5\%$ que mantienen al candidato activo \\
  Rango energético
    & $\Delta E$
    & $H(\max) - H(\min)$ en $[-2,2]^2$ --- contraste entre posiciones extremas \\
  Nodos $A_1$
    & $n_{A_1}$
    & Puntos críticos de silla donde la fibra $H^{-1}(c)$ cambia de topología \\
\end{tabularx}
\end{center}
\medskip

Se dice que $H_i$ \textbf{domina a} $H_j$ si:

\begin{tcolorbox}[codebox]
$H_i \;\geq\; H_j$ \quad en los 4 criterios simultáneamente\\
$H_i \;>\; H_j$   \quad en al menos uno
\end{tcolorbox}

Es decir: $H_i$ es al menos igual de bueno en todo, y estrictamente mejor en
algo. No existe ningún \textit{trade-off} --- $H_i$ gana o empata en cada
dimensión.

% =============================================================================
\section{El algoritmo de capas (Pareto peeling)}

El diagrama se construye por capas, como pelando una cebolla:

\begin{tcolorbox}[codebox]
Frente 1 = \{ H : ningún otro H' domina a H \}\\
\phantom{Frente 1 }$\to$ eliminar Frente 1\\[4pt]
Frente 2 = \{ H restante : ningún otro H' restante domina a H \}\\
\phantom{Frente 2 }$\to$ eliminar Frente 2\\[4pt]
Frente k = repetir hasta vaciar el conjunto
\end{tcolorbox}

Resultado: \textbf{148 frentes} para 296 candidatos. La mayoría de frentes
tienen 1--3 candidatos, lo que indica que el espacio de criterios es altamente
diverso: casi cada Hamiltoniano tiene un perfil único en las 4 dimensiones.

% =============================================================================
\section{Qué es una relación de cobertura (arista del Hasse)}

No se dibuja una arista por cada par donde $H_i$ domina a $H_j$. Solo se dibuja
la \textbf{cobertura directa}: $H_i$ cubre a $H_j$ si:

\begin{tcolorbox}[codebox]
$H_i$ domina a $H_j$\\
Y no existe ningún $H_k$ tal que:\quad $H_i$ dom. $H_k$ \; y \; $H_k$ dom. $H_j$
\end{tcolorbox}

Es decir: el paso de $H_i$ a $H_j$ es un salto irreducible --- no hay
intermediario. Esto da el diagrama más compacto que representa fielmente la
estructura de dominancia.

El diagrama tiene \textbf{475 coberturas} entre 296 candidatos. Las aristas que
cruzan muchos frentes de golpe representan saltos donde no existe ningún
intermediario en ese subespacio de criterios.

% =============================================================================
\section{Cómo leer el diagrama visualmente}

\subsection*{Posición vertical}
\begin{itemize}
  \item \textbf{Arriba} = Frente 1 (los mejores: nadie los domina)
  \item \textbf{Abajo} = Frentes tardíos (dominados por más candidatos)
  \item Cada nivel horizontal es un frente Pareto
\end{itemize}

\subsection*{Posición horizontal}
Dentro de cada frente, los nodos se colocan usando la \textbf{heurística
baricéntrica}: cada nodo se alinea con el promedio horizontal de sus
predecesores. Esto minimiza los cruces de aristas.

\subsection*{Colores de las aristas}

\medskip
\begin{center}
\rowcolors{2}{white}{cTRalt}
\begin{tabular}{L{3cm}C{3cm}L{6cm}}
  \TH{Color} & \TH{Frentes de origen} & \TH{Significado} \\[1pt]
  Oro \texttt{\#FFD700}         & F1       & Coberturas desde los 4 candidatos óptimos \\
  Cian \texttt{\#44DDFF}        & F2       & Segundo nivel del orden \\
  Naranja \texttt{\#FF9944}     & F3       & Tercer nivel \\
  Violeta \texttt{\#CC88FF}     & F4--F8   & Frentes intermedios \\
  Verde menta \texttt{\#88FFAA} & F9--F20  & Zona media \\
  Azul acero \texttt{\#6699CC}  & F21+     & Frentes tardíos \\
\end{tabular}
\end{center}

\subsection*{Los nodos}
\begin{itemize}
  \item \textbf{Versión 2D}: mapa de calor de $H(x,y)$ --- distribución de energía
  \item \textbf{Versión 3D}: superficie $\Gamma(H) = \{z=H(x,y)\}$ ---
        forma topológica
  \item Estrellas naranjas sobre la superficie: nodos $A_1$ (sillas, transiciones)
  \item Borde dorado: Frente 1 \;\textbullet\; Borde cian: Frente 2 \;\textbullet\;
        Borde naranja: Frente 3
\end{itemize}

% =============================================================================
\section{Propiedades matemáticas del poset}

\subsection{No existe supremo ni ínfimo}

El \textbf{supremo} (elemento máximo) sería un candidato que dominara a todos
los demás en los 4 criterios simultáneamente. Ese punto teórico tiene
coordenadas:

\begin{tcolorbox}[codebox]
$\mathrm{sup\;teórico} \;=\; (H_1 = 0.188,\;\; \mathrm{Rob} = 1.0,\;\;
\Delta E = 50.9,\;\; A_1 = 2)$
\end{tcolorbox}

Ningún candidato de los 296 explorados alcanza esos cuatro valores a la vez.
Por tanto \textbf{no existe supremo en el poset}. El \textbf{ínfimo} tampoco
existe: ningún candidato tiene los mínimos simultáneos en las 4 dimensiones.

\subsection{El poset no es una retícula}

Para que el poset fuera una \textbf{retícula} (\textit{lattice}), todo par de
elementos debería tener:
\begin{itemize}
  \item Un \textbf{join} (cota superior mínima): el ``menor'' elemento que
        domina a ambos
  \item Un \textbf{meet} (cota inferior máxima): el ``mayor'' elemento dominado
        por ambos
\end{itemize}

Se verificó computacionalmente que \textbf{ninguno de los 6 pares del Frente 1}
tiene un join dentro del poset. El join existiría solo en la completación
teórica (el supremo externo). Por tanto el poset \textbf{no es retícula}.

\subsection{Teorema de Dilworth: 4 cadenas mínimas}

\begin{tcolorbox}[theorembox]
  \textbf{Teorema de Dilworth} (1950)\medskip\\
  En todo poset finito, el tamaño de la anticadena máxima es igual al mínimo
  número de cadenas necesarias para cubrir el poset.
\end{tcolorbox}

La anticadena máxima es el Frente 1, con \textbf{4 elementos}. Por tanto:

\begin{center}
  {\large\bfseries\color{cNavy} El poset se descompone en exactamente
  \textcolor{cAccent}{4 cadenas}.}
\end{center}

Estas 4 cadenas corresponden a las cuatro columnas aproximadamente verticales
visibles en el diagrama. Cada columna es una secuencia de candidatos donde cada
uno domina al siguiente en al menos un criterio.

\subsection{Clasificación formal}

El poset es un \textbf{poset finito sin cota superior ni inferior}, con:

\medskip
\begin{center}
\rowcolors{2}{white}{cTRalt}
\begin{tabular}{lr}
  \TH{Propiedad} & \TH{Valor} \\[1pt]
  Candidatos (nodos)           & 296 \\
  Frentes Pareto (estratos)    & 148 \\
  Relaciones de cobertura      & 475 \\
  Anticadena máxima (Frente 1) & 4 elementos \\
  Descomposición mínima        & 4 cadenas (Dilworth) \\
  Es retícula                  & No \\
  Tiene supremo                & No \\
  Tiene ínfimo                 & No \\
\end{tabular}
\end{center}

% =============================================================================
\section{El Frente 1: los 4 candidatos óptimos de Pareto}

Los cuatro Hamiltonianos que no son dominados por ningún otro:

\begin{tcolorbox}[codebox]
H\_0094 · score=0.459 · $H_1$=0.188, Rob=1.0, $\Delta E$=31.3, $A_1$=1\\
H\_0113 · score=0.456 · $H_1$=0.173, Rob=1.0, $\Delta E$=35.4, $A_1$=1\\
H\_0045 · score=0.454 · $H_1$=0.166, Rob=1.0, $\Delta E$=43.4, $A_1$=2\\
H\_0042 · score=0.417 · $H_1$=0.120, Rob=1.0, $\Delta E$=33.5, $A_1$=2
\end{tcolorbox}

Son mutuamente incomparables --- forman una \textbf{anticadena}:

\begin{itemize}
  \item H\_0094 tiene el mayor $H_1^{\max}$ pero solo 1 nodo $A_1$
  \item H\_0045 tiene el mayor $\Delta E$ y 2 nodos $A_1$ pero menor $H_1$
  \item H\_0042 tiene 2 nodos $A_1$ (máximo) pero el menor $H_1$ del grupo
  \item H\_0113 es intermedio en todo --- ninguno lo domina
\end{itemize}

Ninguno es ``el mejor'' en todos los sentidos. Son los \textbf{mejores
compromisos posibles} dado el espacio explorado.

% =============================================================================
\section{Qué significa esto para el juego de Go}

\subsection{La energía $H(s_0, s_1)$ como función de evaluación}

En el modelo de Go, el Hamiltoniano $H(s_i, s_j)$ asigna una energía a cada par
de intersecciones vecinas. Hamiltonianos con mayor riqueza topológica
($H_1^{\max}$ alto) pueden distinguir más finamente entre posiciones estratégicas.

\subsection{Significado de cada criterio}

\textbf{$H_1^{\max}$ --- complejidad táctica.}
Mide el tiempo de vida del generador $H_1$ en la filtración de sublevel sets.
Un $H_1^{\max}$ alto indica que la función de energía genera bucles topológicos
persistentes al subir el umbral $c$. En Go: el Hamiltoniano puede detectar
\textbf{estructuras cíclicas de influencia} en el tablero --- grupos que
rodean a otros, cadenas de conexión, patrones de seki.

\textbf{Robustez --- estabilidad del modelo.}
Un candidato robusto (Rob $\approx 1.0$) mantiene su carácter cualitativo ante
pequeñas variaciones de sus coeficientes. Los 4 del Frente 1 tienen Rob = 1.0:
sus propiedades topológicas son \textbf{estructuralmente estables}, no artefactos
de coeficientes precisos.

\textbf{$\Delta E$ --- contraste energético.}
El rango $H(\max) - H(\min)$ en $[-2,2]^2$ mide cuánto separa el Hamiltoniano
las posiciones extremas. En Go: la diferencia energética entre una posición
fuerte (conexión, territorio establecido) y una posición débil (grupos
amenazados, ko). Mayor $\Delta E$ = \textbf{mayor poder discriminativo}.

\textbf{Nodos $A_1$ --- transiciones estratégicas.}
Un nodo $A_1$ es un punto crítico donde $\det(\mathrm{Hess}\,H) < 0$ --- una
silla en la superficie de energía. Al cruzar el valor crítico $c^*$
correspondiente, la fibra $H^{-1}(c)$ cambia de topología (toro $\to$
figura-8). Con 2 nodos $A_1$ (H\_0042, H\_0045), el modelo puede detectar
\textbf{dos transiciones distintas} por partida.

\subsection{Las 4 cadenas de Dilworth y los 4 estilos de juego}

\begin{tcolorbox}[codebox]
Cadena 1 (H\_0094 $\to$ \ldots): linaje con $H_1$ máximo ---
    detecta estructuras cíclicas\\
Cadena 2 (H\_0113 $\to$ \ldots): linaje equilibrado ---
    compromiso entre todos los criterios\\
Cadena 3 (H\_0045 $\to$ \ldots): linaje con $\Delta E$ máximo y 2 $A_1$ ---
    mayor contraste y sensibilidad\\
Cadena 4 (H\_0042 $\to$ \ldots): linaje con $A_1$ múltiple ---
    dos transiciones de régimen
\end{tcolorbox}

Cada cadena representa una \textbf{familia de modelos relacionados por
dominancia}: bajar en una cadena significa perder algo (un criterio se degrada)
pero manteniendo el ``estilo'' general del modelo padre. Los candidatos de
distintas cadenas son \textbf{estratégicamente incomparables}: sirven para
propósitos distintos y no hay motivo matemático para preferir uno sobre otro
sin contexto adicional.

\subsection{Los saltos largos en el diagrama}

Las líneas doradas que cruzan 30--40 frentes de golpe revelan algo importante:
entre H\_0094 (Frente 1) y ciertos candidatos del Frente 40, \textbf{no existe
ningún candidato intermedio} que domine a uno y sea dominado por el otro en los
4 criterios. El espacio de Hamiltonianos cúbicos tiene ``vacíos'' --- regiones
donde no existen compromisos graduales entre ciertos perfiles de calidad.

En términos de Go: no siempre existe una secuencia de modelos de evaluación que
transite suavemente entre ``detecta bien los ciclos'' y ``detecta bien el
contraste energético''. A veces el salto es abrupto.

\subsection{El supremo ausente como objetivo de diseño}

El hecho de que el supremo teórico
$(H_1=0.188,\; \mathrm{Rob}=1.0,\; \Delta E=50.9,\; A_1=2)$
no exista como candidato real es una guía de diseño: \textbf{ese punto es el
Hamiltoniano ideal}. Ningún Hamiltoniano cúbico muestreado lo alcanza, lo que
sugiere que:

\begin{enumerate}
  \item El espacio de búsqueda necesita ampliarse (más muestras, otros templates)
  \item O bien ese punto es matemáticamente imposible con polinomios cúbicos
        de la forma explorada
  \item O se requiere una familia diferente (grado 4, combinaciones no lineales)
\end{enumerate}

El diagrama de Hasse convierte esa imposibilidad en algo visible: el ``vacío''
en la cima del diagrama es el lugar donde debería estar el supremo pero no está.

% =============================================================================
\section{Resumen}

\medskip
\begin{center}
\rowcolors{2}{white}{cTRalt}
\begin{tabularx}{0.96\textwidth}{L{3.2cm}L{4.5cm}X}
  \TH{Concepto} & \TH{Significado matemático} & \TH{Significado en Go} \\[1pt]
  Frente 1
    & Anticadena de elementos maximales del poset
    & Los 4 Hamiltonianos con el mejor compromiso posible \\
  Arista del Hasse
    & Cobertura directa: dominancia sin intermediario
    & Salto irreducible de calidad entre dos modelos \\
  148 frentes
    & Estratificación del poset
    & 148 niveles de calidad estratégica \\
  4 cadenas (Dilworth)
    & Descomposición mínima en cadenas totales
    & 4 ``linajes'' de modelos --- 4 estilos de evaluación \\
  Sin supremo
    & No existe elemento que domine a todos
    & No existe el Hamiltoniano perfecto en el espacio explorado \\
  Sin ínfimo
    & No existe elemento dominado por todos
    & No existe el peor Hamiltoniano absoluto \\
  No es retícula
    & Los joins no existen dentro del poset
    & No siempre hay un ``mejor compromiso'' entre dos modelos incomparables \\
  Nodo $A_1$
    & Punto de silla: $\det(\mathrm{Hess}\,H) < 0$
    & Transición crítica donde la topología del tablero cambia \\
  $H_1^{\max}$ alto
    & Ciclos topológicos persistentes en la fibración
    & Capacidad de detectar estructuras cíclicas de influencia \\
  $\Delta E$ alto
    & Gran rango en $[-2,2]^2$
    & Alto poder discriminativo entre posiciones fuertes y débiles \\
\end{tabularx}
\end{center}

\vfill
\noindent
{\color{cRule}\hrule height 0.5pt}
\vspace{4pt}
{\color{cMuted}\small\itshape
  Pipeline: \texttt{experiments/06\_hamiltonian\_families/} \quad\textbullet\quad
  Figuras: \texttt{output/figures/hasse\_diagram.png} (2D) \;\textbullet\;
           \texttt{hasse\_diagram\_3d.png} (3D)\\
  Catálogo: \texttt{output/catalog.json} \textbullet{} 305 entradas totales
}

\end{document}
